\section{Verification status, limitations, and research program}
\label{sec:status}

A framework of this scope is useful only if its claims are sharply separated.
This section records what the artifact establishes and what remains to be
checked.

\subsection{Established at the human-proof level}

The report gives complete arguments for the following generic conditional
results:
\begin{enumerate}
  \item stabilization and leastness of the finite unified demand analysis;
  \item safe pruning relative to sound transfer summaries;
  \item equivalence of open contracts and ordinary sequents;
  \item acyclic proof-graph compilation by simultaneous cut;
  \item local certificate sealing, stable observation, and bounded named
        export;
  \item fixed bounded-formula reflection from reflected atoms;
  \item soundness of finite positive strategy certificates;
  \item normalized quantifier-free trace reflection;
  \item finite backward viability, least compatible paths, and finite
        observational descent;
  \item guarded quantitative operation interfaces;
  \item the global DARC soundness and relative residual-exactness theorems.
\end{enumerate}
These are theorem schemas: an application must instantiate their stated local
proof obligations.

\subsection{Executable finite evidence}

The reference module \pathcode{reference/darc_reference.py} implements:
\begin{itemize}
  \item finite product-chain demands and threshold rules;
  \item least inflationary closure with a rank bound;
  \item exhaustive monotonicity and leastness checks for small instances;
  \item bounded least search;
  \item positive bounded-formula strategy construction and checking;
  \item deterministic trace execution/checking;
  \item finite backward viability and least paths.
\end{itemize}
The test suite exercises cyclic closure, query-specific liveness, residual
reporting, malformed-rule rejection, strategy tampering, traces, and viable or
nonviable relations.  The generated JSON example records the demand history and
finite certificates.

This code is not an \(\iPRA\) proof checker and is not in the trusted proof
base.  Its role is to test finite algorithms and make the compilation plan
auditable.

\subsection{Lean status}

The structural kernel is written without placeholders and uses current API
names found in the pinned repository.  It was not elaborated here because the
runtime did not contain Lean.  Therefore the artifact does not claim it as a
new checked theorem.  The demand, bounded-reflection, trace, relational, and
compiler files are explicit blueprints containing marked \code{sorry}
milestones.

The theorem ledger in \pathcode{metadata/THEOREM_LEDGER.csv} assigns one of:
fully argued, checked upstream interface, source candidate, Lean blueprint,
executable finite test, or open residual.  No status is inferred from file
names.

\subsection{Deliberate nonclaims}

The artifact does not claim:
\begin{itemize}
  \item a general effective Faltings, Siegel, or Mordell--Weil theorem;
  \item unrestricted exact branch, root, or eigenvector selection;
  \item a global Cauchy modulus from metastability alone;
  \item an infinite path from compatible finite horizons alone;
  \item an effective height bound from qualitative finiteness;
  \item a total universal primitive-recursive evaluator;
  \item that the supplied Lean blueprints compile;
  \item that any new main application theorem is kernel checked.
\end{itemize}
The framework identifies the exact producer interfaces needed by such claims.

\subsection{Research priorities}

The highest-value next step is not another large paper.  It is an integrated
Lean slice that proves one nontrivial endpoint through all layers:
\begin{enumerate}
  \item a fixed bounded formula with alternation;
  \item a loop whose invariant contains the generated truth bit;
  \item a local existential certificate that is sealed;
  \item one named witness selected by checked bounded search;
  \item a cyclic precision demand exposing one residual producer;
  \item a compiled ordinary open sequent;
  \item an open Bridge export to a Foundation client.
\end{enumerate}
Such a pilot would validate the architecture more decisively than separately
elaborating many utility modules.

A second priority is to attach transfer summaries to existing checked
\pathcode{IPRADerivations} lemmas.  The arithmetic library already contains
many proof components; a compact resource interface would make them available
to generated plans without recoding their proofs.

A third priority is producer-driven migration of analytic packets.  Select one
packet whose algebraic and finite parts are mature and whose only missing input
is a clearly stated tail or domain modulus.  Compile the exact conditional
endpoint, then formalize that producer.  This would demonstrate the practical
value of residualization.

\subsection{Conclusion}

DARC turns the broad question ``can this mathematics be expressed in
\(\iPRA\)?'' into a sequence of reviewable obligations:
\begin{enumerate}
  \item What is the exact finite public query?
  \item Which witnesses, precisions, guards, and producers does it force?
  \item Which bounded formulas must be reflected to bits?
  \item Which cycles are finite demand cycles, which are induction loops, and
        which require a genuine extension theorem?
  \item Which local results are ordinary open derivations?
  \item Which deep quantitative inputs remain residual?
\end{enumerate}
The least demand state provides one answer shared by proof scheduling,
certificate escape, precision propagation, and residual exposure.  The final
object is not a semantic promise but an ordinary \(\Prf\;\PRAAx\) derivation.
That combination gives a plausible path from the repository's many rich proof
outlines to a smaller number of reusable, checked, and accurately scoped
formal theorems.
